\documentclass[12pt, a4paper]{article}

% Pacotes essenciais
\usepackage[utf8]{inputenc}
\usepackage[T1]{fontenc}
\usepackage[portuguese]{babel}
\usepackage{amsmath, amssymb}
\usepackage{graphicx}
\usepackage{hyperref}
\usepackage{geometry}

\geometry{margin=2.5cm}

% Informações do documento
\title{Exemplo de Documento \LaTeX}
\author{Seu Nome}
\date{\today}

\begin{document}

\maketitle

\begin{abstract}
Este é um exemplo de documento lorem ipsom \LaTeX{} demonstrando os recursos básicos da linguagem, incluindo formatação de texto, equações matemáticas, listas e referências.
\end{abstract}

\tableofcontents
\newpage

\section{Introdução}

\LaTeX{} é um sistema de preparação de documentos amplamente utilizado na academia para a produção de documentos técnicos e científicos de alta qualidade.

\subsection{Por que usar \LaTeX?}

Existem várias razões para usar \LaTeX:

\begin{itemize}
    \item Excelente qualidade tipográfica
    \item Formatação automática de equações matemáticas
    \item Gerenciamento automático de referências
    \item Consistência visual em todo o documento
\end{itemize}

\section{Formatação de Texto}

Você pode formatar texto de várias maneiras:

\begin{enumerate}
    \item \textbf{Texto em negrito}
    \item \textit{Texto em itálico}
    \item \underline{Texto sublinhado}
    \item \texttt{Texto monoespaçado}
\end{enumerate}

\section{Equações Matemáticas}

\LaTeX{} é especialmente poderoso para equações matemáticas.

\subsection{Equações Inline}

A famosa equação de Einstein é $E = mc^2$.

\subsection{Equações em Bloco}

A fórmula quadrática é dada por:

\begin{equation}
    x = \frac{-b \pm \sqrt{b^2 - 4ac}}{2a}
    \label{eq:quadratica}
\end{equation}

O teorema de Pitágoras:

\begin{equation}
    a^2 + b^2 = c^2
\end{equation}

Uma integral definida:

\begin{equation}
    \int_0^{\infty} e^{-x^2} dx = \frac{\sqrt{\pi}}{2}
\end{equation}

\section{Tabelas}

\begin{table}[h]
    \centering
    \caption{Exemplo de Tabela}
    \begin{tabular}{|c|c|c|}
        \hline
        \textbf{Item} & \textbf{Quantidade} & \textbf{Preço} \\
        \hline
        Produto A & 10 & R\$ 25,00 \\
        Produto B & 5 & R\$ 50,00 \\
        Produto C & 20 & R\$ 15,00 \\
        \hline
    \end{tabular}
    \label{tab:exemplo}
\end{table}

\section{Figuras}

Para incluir figuras, use o ambiente \texttt{figure}:

\begin{verbatim}
\begin{figure}[h]
    \centering
    \includegraphics[width=0.5\textwidth]{imagem.png}
    \caption{Descrição da figura}
    \label{fig:exemplo}
\end{figure}
\end{verbatim}

\section{Conclusão}

Este documento demonstrou os recursos básicos do \LaTeX. Para mais informações, consulte a documentação oficial ou tutoriais online.

\section*{Referências}

\begin{thebibliography}{9}
    \bibitem{lamport94}
    Leslie Lamport,
    \textit{\LaTeX: A Document Preparation System},
    Addison Wesley, 1994.
    
    \bibitem{knuth84}
    Donald E. Knuth,
    \textit{The \TeX book},
    Addison Wesley, 1984.
\end{thebibliography}

\end{document}
